\section{Objetivos}
\label{sec:intro_objetivos}
A partir de la problemática expuesta, el trabajo plantea dos preguntas de investigación que orientan sus objetivos:
\begin{enumerate}[label=\textbf{PI\arabic*.}]
  \item ¿Clasifican los \glspl{vlm} las lesiones de piel, y en particular el melanoma, tan bien como los
  clasificadores entrenados solo con imágenes?
  \item ¿Los reportes que genera un \gls{vlm} contienen el diagnóstico correcto y en qué medida describen la lesión
  que muestra la imagen?
\end{enumerate}

Para responderlas, el objetivo general es desarrollar y evaluar un sistema basado en modelos de aprendizaje
profundo y \glspl{vlm}, adaptados a partir de modelos preentrenados, capaz de asistir en el diagnóstico de lesiones de piel y en la
generación automática de reportes clínicos estructurados a partir de imágenes, con el fin de apoyar el diagnóstico y la documentación de las lesiones. Se analiza el desempeño del sistema en la clasificación, si el reporte contiene el diagnóstico correcto y si puede ejecutarse con recursos computacionales convencionales.

Los objetivos específicos son:
\begin{enumerate}[label=\textbf{OE\arabic*.}]
  \item Diseñar, implementar y evaluar clasificadores de lesiones de piel basados en \glspl{cnn}, \glspl{vit}
  y \glspl{vlm}, generales y especializados en dermatología, para identificar posibles melanomas.
  \item Diseñar, implementar y evaluar un \gls{vlm}, ajustado a partir de un modelo preentrenado, que integre
  información visual y textual para producir un diagnóstico preliminar y el reporte de la lesión.
  \item Evaluar el desempeño del sistema en distintos escenarios, como imágenes de otro conjunto de datos o casos en
  los que las predicciones de los componentes presentan desacuerdo o evidencia insuficiente para afirmar un diagnóstico, analizando su exactitud diagnóstica, la calidad del texto generado y su costo computacional.
\end{enumerate}

La Tabla~\ref{tab:intro_objetivos_preguntas} relaciona cada objetivo específico con la pregunta que contribuye a
responder y con los capítulos donde se lo describe y se lo evalúa.

\begin{table}[H]
  \centering
  \caption{Relación entre los objetivos específicos, las preguntas y los capítulos.}
  \label{tab:intro_objetivos_preguntas}
  \small
  \begin{tabularx}{\textwidth}{@{}>{\raggedright\arraybackslash}p{0.29\textwidth}p{0.10\textwidth}
    >{\raggedright\arraybackslash}X>{\raggedright\arraybackslash}X@{}}
    \toprule
    \textbf{Objetivo específico} & \textbf{Pregunta} & \textbf{Se diseña y describe en} & \textbf{Se evalúa y responde en} \\
    \midrule
    OE1. Clasificadores de lesiones y de melanoma (\gls{cnn}, \gls{vit}, \gls{vlm}) & PI1 &
      Capítulo~\ref{cp:metodos}, Materiales y métodos &
      Capítulo~\ref{cp:res_clasif}, Resultados de clasificación \\
    \filasep
    OE2. \gls{vlm} para generar reportes & PI2 &
      Capítulo~\ref{cp:metodos}, Materiales y métodos &
      Capítulo~\ref{cp:res_gen}, Resultados de generación \\
    \filasep
    OE3. Evaluación en distintos escenarios & PI1 y PI2 &
      Capítulo~\ref{cp:metodos} (protocolo de evaluación) &
      Capítulos~\ref{cp:res_clasif} y~\ref{cp:res_gen} (resultados) y Capítulo~\ref{cp:discusion} (discusión y
      limitaciones) \\
    \bottomrule
  \end{tabularx}
\end{table}
